%\section{Вступ}
\section*{\centering ВСТУП}\addcontentsline{toc}{section}{ВСТУП}

\textbf{Обґрунтування вибору теми дослідження.} Професійно-технічна освіта (ПТО) відіграє ключову роль у підготовці кваліфікованих кадрів для ринку праці та забезпеченні економічного розвитку країн. Моніторинг освітньої діяльності закладів ПТО є необхідною складовою управління якістю освіти, що дозволяє систематично відстежувати та оцінювати ефективність освітніх процесів. Зростання кількості наукових публікацій із цієї тематики потребує систематичного аналізу дослідницького ландшафту для виявлення основних тенденцій, тематичних кластерів та перспективних напрямів досліджень.

Бібліометричний аналіз є потужним інструментом для кількісного дослідження наукової літератури, що дозволяє виявити структуру дослідницького поля, визначити найбільш впливові роботи та ключові слова, а також простежити динаміку розвитку наукових напрямів. Програмне забезпечення VOSViewer є одним із найпоширеніших інструментів для побудови та візуалізації бібліометричних мереж~\cite{vanEck2010}. Водночас використання мови програмування Python з відповідними бібліотеками (NetworkX, scikit-learn, matplotlib) відкриває можливості для відтворення та розширення бібліометричного аналізу, забезпечуючи повну прозорість та відтворюваність результатів.

Порівняння результатів аналізу, отриманих різними інструментами, є важливим кроком для верифікації висновків та оцінки надійності бібліометричних методів. Таке порівняння також створює методичну основу для розробки інформаційно-аналітичної системи моніторингу освітньої діяльності закладів професійної освіти, що є темою бакалаврської кваліфікаційної роботи автора.

\textbf{Об'єкт дослідження}~-- дослідження моніторингу освітньої діяльності закладів професійно-технічної освіти.

\textbf{Предмет дослідження}~-- бібліометричні характеристики та тематична структура наукових досліджень моніторингу ПТО.

\textbf{Мета дослідження}~-- порівняльний бібліометричний аналіз досліджень моніторингу професійно-технічної освіти засобами VOSViewer та Python для виявлення тематичної структури дослідницького поля та оцінки відтворюваності результатів.

Для досягнення поставленої мети було сформульовано такі \textbf{дослідницькі питання}:

\begin{enumerate}
    \item Яка тематична структура досліджень моніторингу ПТО за даними бібліометричного аналізу ключових слів?
    \item Наскільки точно аналітична система на основі Python відтворює результати VOSViewer (частоти термінів, ко-входження, кластеризацію)?
    \item Які основні тенденції та динаміка досліджень моніторингу ПТО спостерігаються у період 1997--2025 років?
\end{enumerate}

\textbf{Використання генеративного штучного інтелекту.} Відповідно до Закону України <<Про академічну доброчесність>>~\cite{AkademichnaDobrochesnist2024} (№~4742-IX), у роботі використано генеративний штучний інтелект Claude (Anthropic) для допомоги в написанні програмного коду аналітичних скриптів мовою Python, генерації чернеток тексту окремих розділів та редагування. Детальний опис використання ШІ наведено в підрозділі~\ref{sec:genai} та Додатку~\ref{app:genai}.

\textbf{Структура та обсяг роботи.} Робота складається зі вступу, трьох розділів, висновків, списку використаних джерел (\total{citenum} найменувань). Робота містить \totalfigures\ рисунків.
Загальний обсяг курсової роботи~-- \pageref{last_page} сторінок.
